\chapter[Python/NumPy 与矩阵乘法四视角]{Python/NumPy 数值编程\\与矩阵乘法的四种视角}\label{ch:numpy-matmul}

\chapEpigraph{程序必须写给人读，只是顺便让机器执行。}{哈罗德·阿贝尔森 \& 杰拉尔德·萨斯曼，《计算机程序的构造和解释》（SICP）}

\begin{chapIntro}{本章导学：前置基础、学习目标与下游坐标}
\textbf{前置基础}：大一 C 语言基础（循环、数组、内存地址概念）与线性代数基本矩阵乘法定义。

\textbf{学习目标}：完成从“逐元素标量循环”到“数组向量化”的思维切换；彻底扫清 NumPy 工程排错中的三大高频雷区（\texttt{view/copy}、\texttt{axis} 规约方向与广播静默维度膨胀）；掌握矩阵乘法的四种互补视角（尤其是列/行线性组合与外积秩 1 求和），并能流畅读写高维批处理张量与 \texttt{np.einsum} 表达式。

\textbf{下游坐标}：本章是本卷第 2--6 章与后续第 2--6 卷所有代码实现、Transformer 注意力计算、SVD/LoRA 低秩分解和神经网络反向传播推导的底层通用语言。
\end{chapIntro}

在 C 语言里，你想把两个 $100 \times 100$ 的矩阵相乘，第一反应是写三层 \texttt{for} 循环，把每个标量元素一个个算出来。但在人工智能与科学计算的世界里，\textbf{矩阵和张量本身就是基本运算单元}：手写 Python 三层循环不仅冗长难读，还会比底层向量化慢上数千倍。本章的目标不是把 Python 语言从头讲一遍，而是站在你已有的 C 语言与基础线代起点上，直接打通 AI 编程与读论文必需的核心表达习惯。

\section{从 C 思维到 Python 工程惯用法}\label{sec:c-to-python}

\begin{recallBox}{C 语言与 Python 的核心心智差异}
C 语言是静态编译型语言：变量必须先声明类型、数组是一段连续裸内存、循环几乎没有解释器开销。Python 是动态解释型语言：变量名只是指向堆上对象的引用，\texttt{for} 循环每步都要做动态类型检查；因此 Python 的黄金法则是\textbf{用简洁的高层语法组织控制流，把密集数值循环全部下沉到 C/NumPy 内核}。
\end{recallBox}

表~\ref{tab:c-vs-py} 浓缩了从 C 过渡到 AI 方向 Python 编程时最关键的六组差异对照。

\begin{table}[htbp]
\centering
\caption{C 语言基础到 AI 方向 Python/NumPy 工程惯用法对照}\label{tab:c-vs-py}
{\small
\begin{tabularx}{\linewidth}{@{}p{24mm}p{52mm}X@{}}
\toprule
\textbf{编程维度} & \textbf{C 语言习惯} & \textbf{Python/NumPy 工程惯用法} \\
\midrule
变量与类型 & 编译期静态声明 \texttt{double x = 1.0;} & 运行期动态绑定，函数签名用类型标注 \texttt{x: float} 提升可读性 \\
数组与切片 & 裸指针 + 长度参数，手动管理越界 & 左闭右开切片 \texttt{a[start:stop:step]}，支持负索引 \texttt{a[-1]} \\
构造序列 & 先 \texttt{malloc} 再 \texttt{for} 循环逐项赋值 & 列表/字典推导式 \texttt{[k*k for k in range(5)]}，一行表达意图 \\
算法封装 & 结构体 \texttt{struct} + 一组外部函数 & 轻量类（\texttt{\_\_init\_\_} / \texttt{fit} / \texttt{predict}），状态与方法统一 \\
防错检查 & \texttt{if (n != m) return -1;} & 关键节点写 \texttt{assert X.shape == (N, D)}，尽早暴露形状错位 \\
随机数生成 & 全局状态 \texttt{srand(42); rand();} & 独立生成器 \texttt{rng = np.random.default\_rng(42)}，隔离实验状态 \\
\bottomrule
\end{tabularx}}
\end{table}

在进入 NumPy 数组之前，只需牢记以下四条贯穿本套七卷本讲义的 Python 工程惯用法：

\begin{enumerate}
  \item \textbf{切片语义与推导式}：切片 \texttt{arr[i:j]} 永远遵循\textbf{左闭右开} $[i, j)$，长度恰好为 $j - i$；负索引 \texttt{-1} 表示最后一个元素，\texttt{arr[::-1]} 表示逆序视图。推导式如 \texttt{\{k: v for k, v in items if v > 0\}} 用于轻量数据整理，但\textbf{切勿用列表推导式替代大规模矩阵运算}。
  \item \textbf{类型标注与形状断言（Shape Assertions）}：AI 代码最常见的 Bug 不是语法报错，而是数组维度静默对错。养成在函数入口写类型标注与 \texttt{assert} 的习惯——例如 \texttt{assert X.ndim == 2}、\texttt{assert X.shape[1] == w.shape[0]}——能在错误扩散到损失函数前当场捕获。
  \item \textbf{统一模型类接口}：从第 2 卷开始，无论是线性回归、逻辑回归还是神经网络模块，我们都沿用 \texttt{fit(X, y)}（拟合参数）与 \texttt{predict(X)}（前向预测）的轻量类封装，便于无缝衔接 scikit-learn 与 PyTorch。
  \item \textbf{现代可复现随机数}：统一调用 \texttt{np.random.default\_rng(seed)} 创建独立生成器 \texttt{rng}，告别旧教材中的全局函数 \texttt{np.random.seed(42)}。这样即使多个函数交替采样，也不会互相污染随机流，保证实验结果百分百可复现。
\end{enumerate}

代码~\ref{lst:py-idiom} 给出了一个 20 行左右的迷你特征标准化器，完整演示了上述四条惯用法。

\begin{lstlisting}[caption={融合类型标注、形状断言、类封装与独立随机生成器的 Python 惯用法}, label={lst:py-idiom}]
import numpy as np

class SimpleScaler:
    """将特征矩阵按列标准化为零均值、单位标准差。"""
    def __init__(self, eps: float = 1e-8) -> None:
        self.eps = eps
        self.mean_: np.ndarray | None = None
        self.std_: np.ndarray | None = None

    def fit(self, X: np.ndarray) -> "SimpleScaler":
        assert X.ndim == 2, f"期望二维矩阵 (N, D)，实际维度 {X.ndim}"
        self.mean_ = X.mean(axis=0, keepdims=True)   # 形状 (1, D)
        self.std_ = X.std(axis=0, keepdims=True)     # 形状 (1, D)
        return self

    def transform(self, X: np.ndarray) -> np.ndarray:
        assert self.mean_ is not None and self.std_ is not None
        assert X.shape[1] == self.mean_.shape[1], "特征维数不匹配"
        return (X - self.mean_) / (self.std_ + self.eps)

rng = np.random.default_rng(42)
X_demo = rng.normal(loc=5.0, scale=2.0, size=(100, 3))
X_norm = SimpleScaler().fit(X_demo).transform(X_demo)
assert np.allclose(X_norm.mean(axis=0), 0.0, atol=1e-7)
\end{lstlisting}

\section{NumPy 核心：\texttt{ndarray} 内存布局、\texttt{axis} 与广播三步规则}\label{sec:numpy-core}

\subsection{\texttt{ndarray} 的底层结构与 \texttt{view/copy} 陷阱}

为什么 Python 原生列表 \texttt{list} 做数值计算很慢？因为 Python 的 \texttt{list} 底层存的并不是连续浮点数，而是一排指向堆内存中分散 \texttt{PyObject} 对象的指针；每次加法都要解引用指针、检查对象类型、拆箱计算再重新装箱。

相比之下，NumPy 的核心数据结构 \texttt{ndarray} 非常贴近 C 语言数组：它在底层维护一块\textbf{连续的同类型原始内存缓冲区（Data Buffer）}，外加一个轻量头部，记录三个核心元信息：
\begin{itemize}
  \item \texttt{shape}（形状）：各轴长度元组，例如 $(N, D)$ 表示 $N$ 个样本、$D$ 个特征；维数记作 \texttt{ndim = len(shape)}。
  \item \texttt{dtype}（元素类型）：决定每个元素占多少字节。科学计算与本卷梯度校验默认使用双精度 \texttt{float64}（8 字节，约 15--17 位有效数字）；深度学习大模型训练常使用单精度 \texttt{float32}（4 字节，约 7 位有效数字）乃至半精度 \texttt{bfloat16} 以减半显存占用。
  \item \texttt{strides}（步幅，单位为字节）：沿每个轴前进 1 步在内存中跳过的字节数。NumPy 默认采用 \textbf{C 连续（行优先，row-major）}布局：对一个形状为 $(N, D)$ 的 \texttt{float64} 数组，第 $(i, j)$ 个元素距离首地址的字节偏移为 $i \cdot (8D) + j \cdot 8$，即 \texttt{strides == (8*D, 8)}。因此按行读取连续命中 CPU 缓存，而跨大步幅按列读取更容易触发缓存未命中。
\end{itemize}

正是因为 \texttt{ndarray} 通过“头部元信息 + 共享数据缓冲区”表示数组，NumPy 在做普通切片和转置（\texttt{X.T}）时\textbf{不会复制底层数据}，而只是返回一个修改了 \texttt{shape}、\texttt{strides} 与起始指针的\textbf{视图（view）}。这带来了零拷贝的高效，但也埋下了初学者最常踩的隐蔽陷阱。

\begin{failBox}{\texttt{view} 与 \texttt{copy} 陷阱：切片修改静默污染原数组}
下面两段看起来非常相似的代码，行为却截然相反：
\begin{lstlisting}[numbers=none]
A = np.array([[1.0, 2.0, 3.0],
              [4.0, 5.0, 6.0]])
s = A[:, 0]          # 基础切片 -> 返回 view（与 A 共享底层内存！）
s[:] = -99.0         # 原地修改 s，A[:, 0] 同时被改成了 [-99., -99.]！

B = np.array([[1.0, 2.0, 3.0],
              [4.0, 5.0, 6.0]])
f = B[[0, 1], 0]     # 花式索引（整数列表/布尔掩码）-> 返回独立 copy
f[:] = -99.0         # B 毫发无损，B[0, 0] 依然是 1.0
\end{lstlisting}
\textbf{排错守则}：（1）基础切片（\texttt{A[1:3, ::2]}）、\texttt{A.T} 与大多数 \texttt{reshape} 返回 \textbf{view}；花式索引（整数数组索引、布尔掩码 \texttt{A[A > 0]}）总是返回 \textbf{copy}。（2）只要你打算原地修改子数组而不影响原数据，\textbf{务必显式调用 \texttt{.copy()}}，不确定时可用 \texttt{np.shares\_memory(a, b)} 直接核实。
\end{failBox}

\subsection{沿轴规约（\texttt{axis}）与 \texttt{keepdims=True}}

在机器学习中，数据表几乎总是按“\textbf{行 = 样本、列 = 特征}”排布成矩阵 $X \in \R^{N \times D}$：第 0 轴（\texttt{axis=0}）沿样本方向（竖着走），第 1 轴（\texttt{axis=1}）沿特征方向（横着走）。理解 \texttt{axis=k} 只有一个绝不记混的口诀：\textbf{\texttt{axis=k} 表示“把第 $k$ 轴压缩（折叠）掉”}。

先看一个 $N=3$ 个样本、$D=2$ 个特征的小例子：
\[
X = \begin{pmatrix} 1 & 10 \\ 3 & 20 \\ 5 & 30 \end{pmatrix} \in \R^{3 \times 2}.
\]
\begin{itemize}
  \item \texttt{X.mean(axis=0)}：压缩第 0 轴（把 3 个样本合为 1 个均值），输出形状从 \texttt{(3, 2)} 变为 \texttt{(2,)}，结果为 \texttt{[3.0, 20.0]}——即\textbf{每个特征的均值}。
  \item \texttt{X.mean(axis=1)}：压缩第 1 轴（把 2 个特征合为 1 个均值），输出形状从 \texttt{(3, 2)} 变为 \texttt{(3,)}，结果为 \texttt{[5.5, 11.5, 17.5]}——即\textbf{每个样本自身特征的均值}。
  \item \texttt{X.mean(axis=1, keepdims=True)}：把第 1 轴长度压成 1 而\textbf{保留轴位}，输出形状为 \texttt{(3, 1)}。这样再写 \texttt{X - X.mean(axis=1, keepdims=True)} 时，\texttt{(3, 2)} 与 \texttt{(3, 1)} 完美对齐；如果不加 \texttt{keepdims=True}，\texttt{(3, 2) - (3,)} 会直接报错！
\end{itemize}

\begin{table}[htbp]
\centering
\caption{数据矩阵 $X \in \R^{N \times D}$ 沿不同 \texttt{axis} 规约的形状与 AI 语义}\label{tab:axis-summary}
{\small
\begin{tabularx}{\linewidth}{@{}p{44mm}p{24mm}X@{}}
\toprule
\textbf{表达式} & \textbf{输出形状} & \textbf{统计与 AI 含义} \\
\midrule
\texttt{X.mean(axis=0)} & \texttt{(D,)} & 跨样本求每个特征的均值（特征标准化） \\
\texttt{X.mean(axis=0, keepdims=True)} & \texttt{(1, D)} & 同上，保留二维形状以便与 $X$ 直接相减 \\
\texttt{X.sum(axis=1, keepdims=True)} & \texttt{(N, 1)} & 对每个样本的所有类别得分求和（如 Softmax 分母） \\
\texttt{X.max(axis=1, keepdims=True)} & \texttt{(N, 1)} & 每个样本的最大 logit（数值稳定 Softmax 减最大值） \\
\texttt{X.mean(axis=None)} & 标量 \texttt{()} & 全部 $N \times D$ 个元素的全局均值 \\
\bottomrule
\end{tabularx}}
\end{table}

\subsection{广播（Broadcasting）三步规则与静默维度灾难}

在实际代码里，我们经常把形状为 \texttt{(100, 3)} 的矩阵加上形状为 \texttt{(3,)} 的偏置向量，或者把 \texttt{(4, 1)} 的列向量加上 \texttt{(1, 5)} 的行向量得到 \texttt{(4, 5)} 的网格。NumPy 允许不同形状的数组做逐元素运算，其背后的严格机制称为\textbf{广播（Broadcasting）}。

\begin{defBox}
\begin{definition}[NumPy 广播三步规则]\label{def:broadcasting}
给定两个参与逐元素运算的数组 $A$ 和 $B$，其形状分别为 $s_A$ 和 $s_B$，NumPy 按以下三步判定并执行广播：
\begin{enumerate}
  \item \textbf{左侧补 1 右对齐}：若两数组维数 \texttt{ndim} 不同，在维数较小的形状元组\textbf{左侧}补 \texttt{1}，直到两者维数相等。
  \item \textbf{逐维兼容性检查}：从左到右逐维比较两形状的轴长 $(d_A, d_B)$。每一维必须满足 $d_A = d_B$，或者其中至少一个为 $1$；否则立即抛出 \texttt{ValueError}。
  \item \textbf{长度为 1 的轴虚拟扩展}：输出数组每一维的长度为 $\max(d_A, d_B)$。凡是原长度为 1 的轴，逻辑上沿该轴重复展开至目标长度——底层实现通过将该轴步幅设为 \texttt{stride = 0} 完成，\textbf{不消耗额外内存复制}。
\end{enumerate}
\end{definition}
\end{defBox}

表~\ref{tab:broadcast-examples} 逐行演练了四组常见形状对齐过程，帮你把三步规则刻进肌肉记忆。

\begin{table}[htbp]
\centering
\caption{广播三步规则的形状对齐演练}\label{tab:broadcast-examples}
{\small
\begin{tabularx}{\linewidth}{@{}p{26mm}p{22mm}p{26mm}p{24mm}X@{}}
\toprule
\textbf{形状 $s_A$} & \textbf{形状 $s_B$} & \textbf{补 1 对齐后} & \textbf{输出形状} & \textbf{判定与典型场景} \\
\midrule
\texttt{(100, 3)} & \texttt{(3,)} & \texttt{(100, 3)} vs \texttt{(1, 3)} & \texttt{(100, 3)} & 合法：给每个样本加特征偏置 \\
\texttt{(4, 1)} & \texttt{(1, 5)} & \texttt{(4, 1)} vs \texttt{(1, 5)} & \texttt{(4, 5)} & 合法：生成两两组合网格/外和 \\
\texttt{(32, 10, 64)} & \texttt{(64,)} & \texttt{(32, 10, 64)} vs \texttt{(1, 1, 64)} & \texttt{(32, 10, 64)} & 合法：Transformer 层偏置广播 \\
\texttt{(100, 3)} & \texttt{(100,)} & \texttt{(100, 3)} vs \texttt{(1, 100)} & \textcolor{ochre}{\textbf{报错}} & 非法：末维 $3 \neq 100$ 且均不为 1 \\
\bottomrule
\end{tabularx}}
\end{table}

广播机制极其优雅，但正因为它“只要满足规则就默默执行”，一旦一维向量 \texttt{(N,)} 与二维列向量 \texttt{(N, 1)} 混用，就会制造出 AI 初学者最致命、最难排查的\textbf{静默 Bug}。

\begin{failBox}{经典静默维度灾难：\texttt{(100, 1) - (100,)} 膨胀成 \texttt{(100, 100)}}
假设你用模型预测了 100 个样本的标签，模型输出 \texttt{y\_pred} 的形状是二维列向量 \texttt{(100, 1)}，而真实标签 \texttt{y\_true} 的形状是一维数组 \texttt{(100,)}。当你写下均方误差：
\begin{lstlisting}[numbers=none]
loss = np.mean((y_pred - y_true) ** 2)   # 看似无害，实则大错特错！
\end{lstlisting}
按定义~\ref{def:broadcasting} 走一遍：\texttt{(100,)} 先左补 1 变成 \texttt{(1, 100)}，再与 \texttt{(100, 1)} 广播——两维都兼容！结果 \texttt{y\_pred - y\_true} 没有报错，而是静默生成了一个 \texttt{(100, 100)} 的成对差值矩阵（内存瞬间膨胀 100 倍），随后 \texttt{np.mean} 把全部 10,000 个跨样本错配误差平均成一个\textbf{完全错误的标量损失}！

\textbf{根治方法}：在计算损失前永远加一句 \texttt{assert y\_pred.shape == y\_true.shape}，或者统一用 \texttt{y\_true.reshape(-1)} 规范为一维向量。
\end{failBox}

\subsection{向量化性能实测：为何 NumPy 比 Python 循环快 5000 倍}

口说无凭，我们在本卷配套实验脚本 \texttt{code/vol1\_experiments.py}（函数 \texttt{exp1\_matmul\_speed}）中，对两个随机双精度矩阵 $A, B \in \R^{120 \times 120}$ 分别运行“纯 Python 三层循环”与“NumPy 矩阵乘法 \texttt{A @ B}”，在本机实测结果如下：

\begin{center}
{\small
\begin{tabular}{@{}lccc@{}}
\toprule
\textbf{实现方式（$120 \times 120$ 双精度矩阵乘）} & \textbf{单次耗时} & \textbf{相对加速比} & \textbf{结果最大绝对偏差} \\
\midrule
纯 Python 三层循环（列表推导式） & $135.7\text{ ms}$ & $1\times$ & \multirow{2}{*}{$1.78 \times 10^{-14}$} \\
NumPy 向量化矩阵乘法 \texttt{A @ B} & $0.027\text{ ms}$ & \textbf{$5{,}021\times$} & \\
\bottomrule
\end{tabular}}
\end{center}

两者计算结果的最大绝对偏差仅为 $1.78 \times 10^{-14}$（源于浮点加法结合律舍入顺序差异，属于双精度机器精度量级），但速度相差\textbf{五千余倍}！从你学过的 C 语言与计算机组成角度看，这 5000 倍加速来自三层红利叠加：
\begin{enumerate}
  \item \textbf{消除解释器开销}：Python 循环执行 $120^3 = 1.728 \times 10^6$ 次乘加时，每次都要做字节码分发、动态类型检查与对象拆装箱；NumPy 一次函数调用直接切入底层编译好的 C/Fortran 例程。
  \item \textbf{缓存分块与连续内存}：底层 BLAS 库的 \texttt{dgemm} 内核将大矩阵拆成贴合 CPU L1/L2 缓存大小的小块计算，大幅减少主存搬运。
  \item \textbf{SIMD 指令与多核并行}：单条 CPU 向量指令（如 AVX/NEON）可在同一时钟周期内同时完成 4--8 组双精度乘加（FMA），大矩阵还会自动启用多核线程并行。
\end{enumerate}

\begin{insightBox}{下游坐标：从 CPU 向量化到 GPU 批处理}
你在本章养成的“消灭 Python 循环、尽量表达成大块连续矩阵运算”的习惯，不仅适用于本卷第 2--6 章与第 2 卷的纯 NumPy 建模，在第 5--6 卷切换到 PyTorch 与 GPU 时依然是同一套底层逻辑——GPU 的 Tensor Core 本质上就是把“大块矩阵乘”做到极致的硬件加速器。
\end{insightBox}

\section{矩阵乘法的四种视角（AI 读论文的母语）}\label{sec:matmul-four-views}

\begin{recallBox}{矩阵乘法的标准定义}
设 $A \in \R^{m \times p}$，$B \in \R^{p \times n}$，则乘积 $C = AB \in \R^{m \times n}$ 的第 $(i, j)$ 个元素定义为
\[
c_{ij} = (AB)_{ij} = \sum_{k=1}^{p} a_{ik} b_{kj}, \qquad i=1,\dots,m,\; j=1,\dots,n.
\]
\end{recallBox}

大一线性代数通常只把上面的公式当作计算口诀来用。然而在阅读人工智能论文和写模型代码时，如果你脑子里只有“第 $i$ 行乘第 $j$ 列”，看到线性层 $Wx$、注意力加权 $PV$、奇异值分解 $\sum \sigma_i u_i v_i^T$ 或反向传播梯度 $\delta x^T$ 时，往往会觉得它们是毫无关系的神秘符号。

事实上，它们只是\textbf{同一个矩阵乘法 $AB$ 的四种等价分解视角}。我们先用图~\ref{fig:matmul-four-views} 一眼看清这四种视角的几何/结构画面，再用同一个 $2 \times 2$ 数值小例逐一验算。

\begin{figure}[htbp]
\centering
\resizebox{\linewidth}{!}{%
\begin{tikzpicture}[
  font=\sffamily\footnotesize,
  mat/.style={draw=ink, thick, fill=white, rounded corners=0.6mm},
  hiA/.style={fill=tealLight, draw=teal, thick},
  hiB/.style={fill=navyLight, draw=navy, thick},
  hiC/.style={fill=ochreLight, draw=ochre, thick},
  panel/.style={draw=tealLine, rounded corners=2mm, fill=white, line width=0.6pt}
]

% ---------- Panel 1: Dot Product ----------
\begin{scope}[shift={(0, 0)}]
  \draw[panel] (0,0) rectangle (7.4, 4.45);
  \node[anchor=north west, font=\sffamily\bfseries\small, text=tealDark] at (0.22, 4.25) {(1) 点积视角：逐元素内积填充};
  % Matrix A
  \draw[mat] (0.45, 1.48) rectangle (1.95, 3.28);
  \draw[hiA] (0.45, 2.48) rectangle (1.95, 2.83);
  \node[text=tealDark] at (1.20, 2.655) {$a_{i\cdot}$};
  \node[below=2pt, text=slate] at (1.20, 1.44) {$A\,(m\times p)$};
  % Matrix B
  \draw[mat] (2.35, 1.48) rectangle (3.85, 3.28);
  \draw[hiB] (3.00, 1.48) rectangle (3.35, 3.28);
  \node[text=navy] at (3.175, 2.38) {$b_{\cdot j}$};
  \node[below=2pt, text=slate] at (3.10, 1.44) {$B\,(p\times n)$};
  % =
  \node[font=\bfseries] at (4.25, 2.38) {$=$};
  % Matrix C
  \draw[mat] (4.65, 1.48) rectangle (6.15, 3.28);
  \draw[hiC] (5.30, 2.48) rectangle (5.65, 2.83);
  \node[text=ochre] at (5.475, 2.655) {\tiny $c_{ij}$};
  \node[below=2pt, text=slate] at (5.40, 1.44) {$C\,(m\times n)$};
  % Formula
  \node[anchor=south, text=ink] at (3.7, 0.16) {$c_{ij} = \inner{a_{i\cdot}}{b_{\cdot j}} = \sum_{k=1}^p a_{ik} b_{kj}$};
\end{scope}

% ---------- Panel 2: Column Linear Combination ----------
\begin{scope}[shift={(7.9, 0)}]
  \draw[panel] (0,0) rectangle (7.4, 4.45);
  \node[anchor=north west, font=\sffamily\bfseries\small, text=tealDark] at (0.22, 4.25) {(2) 列线性组合视角：右乘向量混合列};
  % Matrix A with columns
  \draw[mat] (0.45, 1.48) rectangle (2.25, 3.28);
  \draw[hiA] (0.55, 1.56) rectangle (1.00, 3.20);
  \draw[hiA] (1.13, 1.56) rectangle (1.58, 3.20);
  \draw[hiA] (1.71, 1.56) rectangle (2.16, 3.20);
  \node[text=tealDark] at (0.775, 2.38) {\scriptsize $a_{\cdot 1}$};
  \node[text=tealDark] at (1.355, 2.38) {\scriptsize $a_{\cdot 2}$};
  \node[text=tealDark] at (1.935, 2.38) {\scriptsize $a_{\cdot p}$};
  \node[below=2pt, text=slate] at (1.35, 1.44) {$A$ 的列字典};
  % Vector x
  \draw[hiB] (2.60, 1.68) rectangle (3.10, 3.08);
  \node[text=navy] at (2.85, 2.78) {\scriptsize $x_1$};
  \node[text=navy] at (2.85, 2.38) {\scriptsize $\vdots$};
  \node[text=navy] at (2.85, 1.95) {\scriptsize $x_p$};
  \node[below=2pt, text=slate] at (2.85, 1.44) {$x$};
  % =
  \node[font=\bfseries] at (3.48, 2.38) {$=$};
  % Output
  \node[anchor=west, text=ink] at (3.78, 2.38) {$x_1 \textcolor{teal}{a_{\cdot 1}} + \dots + x_p \textcolor{teal}{a_{\cdot p}}$};
  \node[anchor=south, text=ink] at (3.7, 0.16) {$Ax \in \col(A)$；$AB$ 的第 $j$ 列 $= A b_{\cdot j}$};
\end{scope}

% ---------- Panel 3: Row Linear Combination ----------
\begin{scope}[shift={(0, -4.90)}]
  \draw[panel] (0,0) rectangle (7.4, 4.45);
  \node[anchor=north west, font=\sffamily\bfseries\small, text=tealDark] at (0.22, 4.25) {(3) 行线性组合视角：左乘向量加权行};
  % Row vector x^T
  \draw[hiB] (0.40, 2.18) rectangle (1.95, 2.63);
  \node[text=navy] at (1.175, 2.405) {\scriptsize $x_1 \;\cdots\; x_m$};
  \node[below=2pt, text=slate] at (1.175, 2.02) {$x^T$};
  % Matrix B with rows
  \draw[mat] (2.25, 1.48) rectangle (3.95, 3.28);
  \draw[hiA] (2.33, 2.77) rectangle (3.87, 3.19);
  \draw[hiA] (2.33, 2.17) rectangle (3.87, 2.59);
  \draw[hiA] (2.33, 1.57) rectangle (3.87, 1.99);
  \node[text=tealDark] at (3.10, 2.98) {\scriptsize $b_{1\cdot}$};
  \node[text=tealDark] at (3.10, 2.38) {\scriptsize $b_{2\cdot}$};
  \node[text=tealDark] at (3.10, 1.78) {\scriptsize $b_{m\cdot}$};
  \node[below=2pt, text=slate] at (3.10, 1.44) {$B$ 的行样本};
  % =
  \node[font=\bfseries] at (4.30, 2.38) {$=$};
  % Output
  \node[anchor=west, text=ink, font=\scriptsize] at (4.49, 2.38) {$x_1 \textcolor{teal}{b_{1\cdot}} + \dots + x_m \textcolor{teal}{b_{m\cdot}}$};
  \node[anchor=south, text=ink] at (3.7, 0.16) {$x^T B \in \row(B)$（Attention 权重加权 Value 行）};
\end{scope}

% ---------- Panel 4: Outer Product Sum ----------
\begin{scope}[shift={(7.9, -4.90)}]
  \draw[panel] (0,0) rectangle (7.4, 4.45);
  \node[anchor=north west, font=\sffamily\bfseries\small, text=tealDark] at (0.22, 4.25) {(4) 外积秩 1 求和视角：低秩结构分解};
  % Sum of rank-1
  \node[anchor=west, text=ink] at (0.35, 2.38) {$AB = \displaystyle\sum_{k=1}^{p}$};
  % Col vector
  \draw[hiA] (1.95, 1.58) rectangle (2.40, 3.18);
  \node[text=tealDark] at (2.175, 2.38) {\scriptsize $a_{\cdot k}$};
  % Row vector
  \draw[hiB] (2.55, 2.78) rectangle (4.05, 3.18);
  \node[text=navy] at (3.30, 2.98) {\scriptsize $b_{k\cdot}$};
  % Equals rank-1 block
  \node[font=\bfseries] at (4.42, 2.38) {$=$};
  \draw[hiC] (4.78, 1.63) rectangle (6.12, 3.13);
  \node[align=center, text=ochre] at (5.45, 2.38) {\scriptsize 秩 $\le 1$\\\scriptsize 矩阵块};
  \node[right=1pt, text=ink] at (6.16, 2.38) {$+\cdots$};
  \node[anchor=south, text=ink] at (3.7, 0.16) {$AB = \sum_{k=1}^p a_{\cdot k} b_{k\cdot}$（SVD、LoRA、反传外积梯度）};
\end{scope}

\end{tikzpicture}}
\caption{矩阵乘法的四种互补视角：同一种运算，在 AI 中对应四种截然不同的建模语义}\label{fig:matmul-four-views}
\end{figure}

下面我们用同一对 $2 \times 2$ 矩阵手算验证这四种视角，记
\[
A = \begin{pmatrix} 2 & 1 \\ 0 & 3 \end{pmatrix}, \qquad
B = \begin{pmatrix} 1 & 4 \\ 2 & -1 \end{pmatrix}.
\]

\begin{enumerate}
  \item \textbf{视角一：点积视角（逐元素内积）}。第 $(i,j)$ 个元素等于 $A$ 的第 $i$ 行与 $B$ 的第 $j$ 列做点积：
  \[
  AB = \begin{pmatrix}
    2\cdot 1 + 1\cdot 2 & 2\cdot 4 + 1\cdot(-1) \\
    0\cdot 1 + 3\cdot 2 & 0\cdot 4 + 3\cdot(-1)
  \end{pmatrix}
  = \begin{pmatrix} 4 & 7 \\ 6 & -3 \end{pmatrix}.
  \]
  在 AI 中，这一视角用于理解\textbf{相似度打分矩阵}：若 $Q$ 的每一行是一个查询向量、$K$ 的每一行是一个键向量，则 $QK^T$ 的第 $(i,j)$ 个元素就是第 $i$ 个查询与第 $j$ 个键的点积相似度。

  \item \textbf{视角二：列线性组合视角（右乘列向量混合 $A$ 的列）}。把 $A$ 写成列向量组 $A = [a_{\cdot 1},\, a_{\cdot 2}]$，则 $AB$ 的第 1 列和第 2 列分别是 $A$ 各列的线性组合，组合系数由 $B$ 对应列给出：
  \[
  (AB)_{\cdot 1} = 1\begin{pmatrix}2\\0\end{pmatrix} + 2\begin{pmatrix}1\\3\end{pmatrix} = \begin{pmatrix}4\\6\end{pmatrix}, \qquad
  (AB)_{\cdot 2} = 4\begin{pmatrix}2\\0\end{pmatrix} + (-1)\begin{pmatrix}1\\3\end{pmatrix} = \begin{pmatrix}7\\-3\end{pmatrix}.
  \]
  这一视角直接告诉你一个深刻事实：\textbf{对任意向量 $x$，$Ax$ 永远落在 $A$ 的列空间 $\col(A)$ 中}。第 2 章做线性最小二乘 $\min_w \|Aw - b\|_2^2$ 时，正是靠这一视角把它看成“在列空间 $\col(A)$ 中找离观测向量 $b$ 最近的点”！

  \item \textbf{视角三：行线性组合视角（左乘行向量加权 $B$ 的行）}。把 $B$ 写成行向量组 $B = \begin{psmallmatrix} b_{1\cdot} \\ b_{2\cdot} \end{psmallmatrix}$，则 $AB$ 的第 1 行和第 2 行分别是 $B$ 各行的线性组合，权重由 $A$ 对应行给出：
  \[
  (AB)_{1\cdot} = 2\begin{pmatrix}1 & 4\end{pmatrix} + 1\begin{pmatrix}2 & -1\end{pmatrix} = \begin{pmatrix}4 & 7\end{pmatrix}, \qquad
  (AB)_{2\cdot} = 0\begin{pmatrix}1 & 4\end{pmatrix} + 3\begin{pmatrix}2 & -1\end{pmatrix} = \begin{pmatrix}6 & -3\end{pmatrix}.
  \]
  在 Transformer 自注意力机制 $\operatorname{Attention}(Q,K,V) = P V$ 中（第 5 卷），$V \in \R^{T \times d_v}$ 的每一行是一个词元的值向量，$P \in \R^{T \times T}$ 的第 $i$ 行是经过 Softmax 归一化的注意力概率分布。按行组合视角看，\textbf{输出的第 $i$ 行就是用第 $i$ 行的注意力权重对所有词元值向量做加权平均}！

  \item \textbf{视角四：外积秩 1 求和视角（列乘行分解）}。把乘积拆成 $A$ 的第 $k$ 列与 $B$ 的第 $k$ 行的\textbf{外积（Outer Product）}之和：
  \[
  AB = a_{\cdot 1} b_{1\cdot} + a_{\cdot 2} b_{2\cdot}
  = \begin{pmatrix}2\\0\end{pmatrix}\begin{pmatrix}1 & 4\end{pmatrix}
  + \begin{pmatrix}1\\3\end{pmatrix}\begin{pmatrix}2 & -1\end{pmatrix}
  = \begin{pmatrix}2 & 8\\0 & 0\end{pmatrix}
  + \begin{pmatrix}2 & -1\\6 & -3\end{pmatrix}
  = \begin{pmatrix}4 & 7\\6 & -3\end{pmatrix}.
  \]
  注意观察等式右边的每一项 $u v^T$（$u \in \R^m, v \in \R^n$）：它的每一列都是 $u$ 的标量倍，每一行都是 $v^T$ 的标量倍，因此任何非零外积 $uv^T$ 的秩严格等于 $1$。由此顺带看清一个重要结论：$m \times p$ 矩阵乘以 $p \times n$ 矩阵，本质上是 $p$ 个秩 1 矩阵的叠加，故 $\rank(AB) \le p$。
\end{enumerate}

\begin{insightBox}{为什么外积秩 1 视角是 AI 最值钱的视角？}
至少有三大核心算法直接建立在外积视角之上：
\begin{itemize}
  \item \textbf{奇异值分解与图像压缩（本卷第 4 章）}：把大矩阵写成 $A = \sum_{i=1}^r \sigma_i u_i v_i^T$，只保留前 $k$ 个最大奇异值对应的秩 1 块即可高精度压缩。
  \item \textbf{反向传播梯度（本卷第 5 章与第 5 卷）}：线性层 $y = Wx$ 的权重梯度恰好是外积 $\frac{\partial L}{\partial W} = \delta x^T$；整个小批次 $B$ 个样本的总梯度则是 $\Delta^T X = \sum_{n=1}^B \delta_n x_n^T$，一步矩阵乘汇总所有样本的外积！
  \item \textbf{大模型 LoRA 微调（第 6 卷）}：不直接更新巨大的 $d \times d$ 权重增量 $\Delta W$，而是参数化为低秩乘积 $\Delta W = B A$（$B \in \R^{d \times r},\, A \in \R^{r \times d},\, r \ll d$），本质上是用 $r$ 个秩 1 块逼近参数更新，可训练参数量从 $d^2$ 骤降至 $2dr$。
\end{itemize}
\end{insightBox}

除了上述标准矩阵乘法外，AI 中还频繁使用\textbf{哈达玛积（Hadamard Product，逐元素乘积）}，记作 $A \had B$。对同形状矩阵 $A, B \in \R^{m \times n}$，定义
\[
(A \had B)_{ij} = a_{ij} b_{ij}.
\]
它在深度学习中无处不在：Dropout 掩码 $h \had m$（把随机置 0 的神经元屏蔽掉）、现代大模型（如 LLaMA）的 SwiGLU 门控层 $\operatorname{SiLU}(x W_1) \had (x W_2)$、以及逐元素特征缩放。请牢记 NumPy 运算符的铁律：\textbf{\texttt{A * B} 是逐元素哈达玛积 $A \had B$，而 \texttt{A @ B} 才是矩阵乘法 $AB$}，二者绝不可混淆！

\section{批处理高维张量乘、\texttt{np.einsum} 与基础可视化}\label{sec:batch-einsum}

\subsection{批处理高维张量与 \texttt{@} 的高维规则}

真实 AI 任务很少一次只算单个样本或单个注意力头。为了喂饱现代硬件，我们会把样本批量、时间步、注意力头打包成三维或四维张量（多维 \texttt{ndarray}）：
\begin{itemize}
  \item \textbf{序列批处理张量} \texttt{(B, T, D)}：$B$ 个样本（Batch）、每个样本 $T$ 个词元（Tokens）、每个词元用 $D$ 维向量表示。
  \item \textbf{多头注意力张量} \texttt{(B, H, T, D\_h)}：把 $D$ 维拆成 $H$ 个独立注意力头（Heads），每头维度 $D_h = D/H$。
\end{itemize}

面对高维张量，NumPy 的 \texttt{A @ B} 遵循一条异常清晰的规则：\textbf{把最后两维 \texttt{(-2, -1)} 视作待相乘的二维矩阵，前面的所有批次维度 \texttt{[:-2]} 按广播三步规则（定义~\ref{def:broadcasting}）对齐}。

看两个深度学习中最常见的例子：
\begin{enumerate}
  \item \textbf{批量线性映射}：输入 \texttt{X} 形状为 \texttt{(B, T, D)}，共享权重 \texttt{W} 形状为 \texttt{(D, M)}。执行 \texttt{X @ W} 时，\texttt{W} 的前置维度自动广播为 \texttt{(1, D, M)}，与每个样本的 \texttt{(T, D)} 矩阵相乘，得到 \texttt{(B, T, M)}——等价于对每个样本中的每个词元独立施加同一线性变换 $W$。
  \item \textbf{多头注意力打分矩阵}：查询 \texttt{Q} 形状为 \texttt{(B, H, T, D\_h)}，键 \texttt{K} 形状为 \texttt{(B, H, S, D\_h)}。先令 \texttt{Kt = K.swapaxes(-1, -2)} 把最后两维转置成 \texttt{(B, H, D\_h, S)}，再计算 \texttt{Q @ Kt}，便一次性并行算完所有 $B \times H$ 组注意力打分矩阵，输出形状为 \texttt{(B, H, T, S)}！
\end{enumerate}

\subsection{爱因斯坦求和约定与 \texttt{np.einsum}}

当张量超过三维、涉及多处转置和求和时，连续写 \texttt{swapaxes} 和 \texttt{@} 往往让人眼花缭乱。NumPy 提供的 \texttt{np.einsum}（爱因斯坦求和约定）允许你\textbf{直接把下标字母写出来}，代码几乎就是数学公式本身。

阅读与书写 \texttt{np.einsum("输入下标->输出下标", ...)} 只需掌握两条规则：
\begin{enumerate}
  \item \textbf{出现在箭头左边、但未出现在箭头右边的字母}：表示沿该下标求和（缩并，contraction）。
  \item \textbf{同时出现在箭头左右两边的字母}：表示保留在输出张量中，并按箭头右边的顺序排列轴位。
\end{enumerate}

表~\ref{tab:einsum-cheatsheet} 整理了从本卷到第 6 卷最高频的六种 \texttt{np.einsum} 写法。

\begin{table}[htbp]
\centering
\caption{AI 高频张量运算的 \texttt{np.einsum} 速查表}\label{tab:einsum-cheatsheet}
{\small
\begin{tabularx}{\linewidth}{@{}p{28mm}p{38mm}p{32mm}X@{}}
\toprule
\textbf{运算名称} & \textbf{\texttt{np.einsum} 写法} & \textbf{数学表达式} & \textbf{典型用途与说明} \\
\midrule
向量点积 & \texttt{'i,i->', u, v} & $\sum_i u_i v_i$ & 内积、余弦相似度分子 \\
向量外积 & \texttt{'i,j->ij', u, v} & $(uv^T)_{ij} = u_i v_j$ & 秩 1 矩阵构造、单样本梯度 \\
标准矩阵乘 & \texttt{'ik,kj->ij', A, B} & $\sum_k A_{ik} B_{kj}$ & 等价于 \texttt{A @ B} \\
Frobenius 内积 & \texttt{'ij,ij->', A, B} & $\sum_{i,j} A_{ij} B_{ij}$ & 等于 $\tr(A^T B)$（第 5 章矩阵微分） \\
批量二次型 & \texttt{'ni,ij,nj->n', X, Q, X} & $x_n^T Q x_n$ & 一次性计算 $N$ 个样本的二次损失 \\
多头注意力打分 & \texttt{'bhtd,bhsd->bhts', Q, K} & $\sum_d Q_{b,h,t,d} K_{b,h,s,d}$ & 无需手动 \texttt{swapaxes}，可读性极佳 \\
\bottomrule
\end{tabularx}}
\end{table}

\begin{remark}[多张量缩并的性能提示]
当 \texttt{np.einsum} 同时缩并三个或更多张量时（例如批量二次型表达式 \texttt{'ni,ij,nj->n'}），默认从左到右逐对计算可能产生不必要的巨大中间张量。养成加上参数 \texttt{optimize=True} 的习惯，NumPy 会自动搜索最优加括号顺序，大幅降低计算量与内存峰值。
\end{remark}

\subsection{完整可运行代码范例与基础可视化模板}

代码~\ref{lst:ch1-demo} 整合了本章全部核心技法：校验矩阵乘法四视角等价、防静默广播残差检查、对比 \texttt{@} 与 \texttt{np.einsum} 计算多头注意力打分，并给出本系列通用的最小 \texttt{matplotlib} 绘图骨架。

\begin{lstlisting}[caption={矩阵乘法四视角、多头注意力 \texttt{einsum} 验证与最小可视化模板}, label={lst:ch1-demo}]
import numpy as np
import matplotlib.pyplot as plt

# 1. 矩阵乘法四视角数值一致性校验
A = np.array([[2.0, 1.0], [0.0, 3.0]])
B = np.array([[1.0, 4.0], [2.0, -1.0]])
C_dot = A @ B
C_outer = sum(np.outer(A[:, k], B[k, :]) for k in range(A.shape[1]))
assert np.allclose(C_dot, C_outer, atol=1e-12)

# 2. 防止静默广播的安全均方误差函数
def safe_mse(y_pred: np.ndarray, y_true: np.ndarray) -> float:
    assert y_pred.shape == y_true.shape, (
        f"形状不一致将触发静默广播：{y_pred.shape} vs {y_true.shape}"
    )
    return float(np.mean((y_pred - y_true) ** 2))

# 3. 批处理多头注意力打分：@ 与 np.einsum 完全等价
rng = np.random.default_rng(2026)
B_sz, H, T, S, Dh = 2, 4, 8, 6, 16
Q = rng.normal(size=(B_sz, H, T, Dh))
K = rng.normal(size=(B_sz, H, S, Dh))
scores_at = (Q @ K.swapaxes(-1, -2)) / np.sqrt(Dh)
scores_ein = np.einsum("bhtd,bhsd->bhts", Q, K) / np.sqrt(Dh)
assert scores_at.shape == (B_sz, H, T, S)
assert np.allclose(scores_at, scores_ein, atol=1e-12)

# 4. 最小科研绘图模板（面向对象 Axes 接口）
x = np.linspace(0.0, 2.0 * np.pi, 200)
fig, ax = plt.subplots(figsize=(5.5, 3.2), dpi=150)
ax.plot(x, np.sin(x), color="#0F5C5B", linewidth=2.0, label=r"$\sin(x)$")
ax.plot(x, np.cos(x), color="#9A4624", linestyle="--", label=r"$\cos(x)$")
ax.set_xlabel("x"); ax.set_ylabel("y"); ax.grid(alpha=0.25); ax.legend()
fig.tight_layout()
# fig.savefig("trig_demo.png")
\end{lstlisting}

\begin{extBox}{步幅 \texttt{strides} 零拷贝滑窗直觉（\texttt{sliding\_window\_view}）}
前面提到，广播底层靠将步幅 \texttt{stride} 设为 0 实现零内存拷贝。更进一步，NumPy 提供了 \texttt{np.lib.stride\_tricks.sliding\_window\_view(x, window\_shape)}：对一维时间序列或二维图像，只需巧妙调整头部 \texttt{shape} 与 \texttt{strides}，就能零拷贝生成所有滑动窗口视图——这正是第 5 卷卷积神经网络（CNN）底层 \texttt{im2col} 向量化加速的思想雏形。初读只需知道其存在，二刷或写手写卷积时再深入尝试即可。
\end{extBox}

\section{本章习题}\label{sec:ch01-exercises}

\exTier{A. 概念辨析与反例查错}{先独立思考并定位陷阱，切忌直接翻答案}

\begin{enumerate}[label=\textbf{1.\arabic*}]
  \item \textbf{广播规则与静默维度错误诊断}：
  （1）根据广播三步规则，逐一判断以下四组形状能否合法广播；若合法，写出输出形状，若非法，说明在哪一维冲突：
  \begin{center}
  (a) \texttt{(8, 1, 5)} 与 \texttt{(4, 5)}；\quad
  (b) \texttt{(8, 3, 5)} 与 \texttt{(3, 1)}；\quad
  (c) \texttt{(10, 4)} 与 \texttt{(10,)}；\quad
  (d) \texttt{(64, 1)} 与 \texttt{(64,)}。
  \end{center}
  （2）设某回归模型在 64 个样本上的预测值为 \texttt{y\_pred}（形状 \texttt{(64, 1)}），真实标签为 \texttt{y\_true}（形状 \texttt{(64,)}）。解释直接执行 \texttt{np.mean((y\_pred - y\_true)**2)} 会发生什么数值错误，并给出两种安全的修复写法。

  \item \textbf{\texttt{view/copy} 与矩阵乘法运算符排错}：阅读下面这段试图实现“归一化副本并计算矩阵乘积 $AB$”的 6 行代码，在不运行机器的前提下，逐行写出 \texttt{A[0, 0]} 的最终值与 \texttt{C} 的计算结果，指出其中隐藏的\textbf{两处 Bug} 并改正：
  \begin{lstlisting}[numbers=none]
A = np.array([[2.0, 4.0], [1.0, 3.0]])
B = np.array([[1.0, 0.0], [0.0, 2.0]])
A_scaled = A[:]          # Bug 1?
A_scaled /= 2.0
C = A_scaled * B         # Bug 2?
print(A[0, 0], C)
  \end{lstlisting}
\end{enumerate}

\exTier{B. 手算小例与基础演练}{亲手算一遍，把概念变成直觉}

\begin{enumerate}[label=\textbf{1.\arabic*}, resume]
  \item \textbf{矩阵乘法四视角手算验证}：设
  \[
  A = \begin{pmatrix} 1 & 2 \\ 0 & -1 \\ 3 & 1 \end{pmatrix}, \qquad
  B = \begin{pmatrix} 2 & -1 \\ 1 & 3 \end{pmatrix}.
  \]
  请分别使用（1）\textbf{列线性组合视角}、（2）\textbf{行线性组合视角}和（3）\textbf{外积秩 1 求和视角}计算乘积 $AB$，验证三者结果完全一致，并说明外积展开中每一项的矩阵秩是多少。

  \item \textbf{\texttt{axis} 规约与 \texttt{np.einsum} 翻译}：
  （1）设数据矩阵 $X = \begin{psmallmatrix} 2 & 10 & 4 \\ 6 & 20 & 4 \end{psmallmatrix}$，手算 \texttt{X.mean(axis=0, keepdims=True)}、\texttt{X.std(axis=0, keepdims=True)} 以及按列中心化结果 \texttt{X - X.mean(axis=0, keepdims=True)}。
  （2）将以下三个 \texttt{np.einsum} 表达式翻译为等价的数学求和式和常规 NumPy 写法（使用 \texttt{@}、\texttt{*} 或 \texttt{sum}），并写出输出形状（设 $A:(m,n)$，$x:(n,)$，$U,V:(N,D)$，$H:(B,T,D)$，$W:(D,M)$）：
  \begin{center}
  (a) \texttt{np.einsum('ij,j->i', A, x)}；\quad
  (b) \texttt{np.einsum('ni,ni->n', U, V)}；\quad
  (c) \texttt{np.einsum('btd,dm->btm', H, W)}。
  \end{center}
\end{enumerate}

\exTier{C. 实用证明与理论推导}{紧扣 AI 核心性质，无冷门偏题}

\begin{enumerate}[label=\textbf{1.\arabic*}, resume]
  \item \textbf{秩 1 外积矩阵的基本性质}：设非零向量 $u \in \R^m$，$v \in \R^n$，令 $M = u v^T \in \R^{m \times n}$。证明：
  （1）$\rank(M) = 1$；
  （2）当 $m = n$ 时，矩阵的迹满足 $\tr(u v^T) = u^T v$；
  （3）对任意矩阵 $A \in \R^{m \times p}$ 与 $B \in \R^{p \times n}$，利用外积求和视角证明 $\rank(AB) \le p$。
  \exHint{第（1）问观察 $M$ 的每一列均为 $u$ 的倍数且至少有一列非零；第（2）问直接展开对角元 $M_{ii} = u_i v_i$ 求和；第（3）问利用秩的次可加性 $\rank(X + Y) \le \rank(X) + \rank(Y)$。}

  \item \textbf{低秩分解的乘法加括号顺序与 LoRA 计算量}：设大型权重更新矩阵采用低秩参数化 $\Delta W = B A$，其中 $B \in \R^{d \times r}$、$A \in \R^{r \times d}$ 且 $r \ll d$。给定单个输入向量 $x \in \R^d$，比较两种计算顺序：方案一先显式合成 $M = BA$ 再算 $Mx$；方案二先算 $z = Ax$ 再算 $Bz$。分别计算两种方案所需的标量乘法次数，并在 $d = 4096,\, r = 16$ 时给出具体加速比。
  \exHint{一个 $p \times q$ 矩阵乘以 $q \times s$ 矩阵恰好需要 $p q s$ 次标量乘法；注意向量可看作 $d \times 1$ 矩阵。}
\end{enumerate}

\exTier{D. 数值编程与可自检实验}{独立跑通并让全部 \texttt{assert} 通过}

\begin{enumerate}[label=\textbf{1.\arabic*}, resume]
  \item \textbf{零循环实现成对欧氏距离矩阵}：给定两组点集 $X \in \R^{N \times D}$ 与 $Y \in \R^{M \times D}$，要求计算成对欧氏距离矩阵 $D \in \R^{N \times M}$，其中 $D_{ij} = \|x_i - y_j\|_2$。请分别实现以下两种零循环方法：
  \begin{itemize}
    \item 方法一（三维广播法）：通过 \texttt{X[:, None, :] - Y[None, :, :]} 直接广播求差再沿末轴求范数；
    \item 方法二（展开恒等式法）：利用 $\|x_i - y_j\|_2^2 = \|x_i\|_2^2 + \|y_j\|_2^2 - 2 x_i^T y_j$，借助矩阵乘法 \texttt{X @ Y.T} 与广播一次性算出平方距离，再用 \texttt{np.maximum(..., 0.0)} 截断浮点舍入产生的微小负数后开方。
  \end{itemize}
  \textbf{自检标准}：使用 \texttt{rng = np.random.default\_rng(0)} 生成 $N=120,\, M=80,\, D=32$ 的标准正态矩阵，将两种方法与双重循环基准结果对比，验证最大绝对误差均满足 \texttt{max\_abs\_err < 1e-12}，并比较两种方法在 $N=M=1000,\, D=64$ 时的耗时差异。
  \exHint{方法二中令 \texttt{x2 = np.sum(X**2, axis=1, keepdims=True)}（形状 \texttt{(N, 1)}）与 \texttt{y2 = np.sum(Y**2, axis=1)[None, :]}（形状 \texttt{(1, M)}），即可安全广播相加。}

  \item \textbf{多头缩放点积注意力的纯 NumPy 实现与校验}：编写函数 \texttt{multihead\_attention(Q, K, V)}，其中输入 \texttt{Q, K, V} 形状均为 \texttt{(B, H, T, D\_h)}，按以下公式计算输出 $O \in \R^{B \times H \times T \times D_h}$：
  \[
  S = \frac{Q K^T}{\sqrt{D_h}}, \qquad
  P = \operatorname{softmax}_{\text{行}}(S), \qquad
  O = P V.
  \]
  要求：（1）行 Softmax 必须先减去每行最大值（\texttt{keepdims=True}）防止 \texttt{np.exp} 溢出；（2）分别用高维 \texttt{@} 与 \texttt{np.einsum} 各实现一次。
  \textbf{自检标准}：使用 \texttt{rng = np.random.default\_rng(7)} 在 $B=2,\, H=4,\, T=16,\, D_h=8$ 上测试，验证：（a）输出形状严格为 \texttt{(2, 4, 16, 8)}；（b）两种实现的最大绝对误差 \texttt{< 1e-12}；（c）注意力权重矩阵 $P$ 每一行之和与 $1$ 的最大偏差 \texttt{< 1e-12}，且将 $S$ 整体加上常数 $1000$ 后 $P$ 保持不变。
  \exHint{利用 \texttt{np.einsum} 时，打分步用 \texttt{'bhtd,bhsd->bhts'}，加权求和步用 \texttt{'bhts,bhsd->bhtd'}。}
\end{enumerate}
